\chapter{Introduction}

\section{Purpose and Goal of Project}

The purpose of this project was to design, build and test a low-cost rover prototype capable of remote operation and mapping on a simulated lunar surface. 
The rover features a differential drive system, a rocker-bogie-inspired suspension, closed-loop stepper motors and a LiDAR mapping system.

A key mapping feature of the rover is the ability to rotate a 2D LiDAR sensor to create a pseudo-3D point cloud of its surroundings. 
The goal of this project was to investigate whether a relatively simple and 
low-cost rover prototype could support remote-controlled navigation in a simulated lunar test context and generate usable pseudo-3D point-cloud data in a 
low-light environment.

\section{Report Structure}
This report is structured according to the thesis guidelines provided on Canvas. Chapter 2 presents the background, problem statement, objectives and limitations.
Chapter 3 describes the relevant theory behind our project, LiDAR mapping, motor control, communications, simulation, power systems and finite element analysis.
Chapter 4 presents the materials and methods used to design, build and test the rover.
Chapter 5 presents the results from subsystem testing, simulation, physical assembly and mechanical analysis.
Chapter 6 discusses the results, limitations and possible sources of error.
Chapter 7 concludes the report and presents recommendations for further work.
